% F - three baselines side by side (Chapter 2): MLP vs GRU vs SNN, all consuming
% the SAME length-T multivariate input so they are directly comparable.
% MLP: per-timestep MLP applied independently each step (NO hidden state) ->
%      mean over time -> logits.
% GRU: unrolled recurrent cell with real-valued hidden state h_t -> logits.
% SNN: unrolled LIF (lifdot) emitting spikes with beta state -> mean over time -> logits.
\documentclass[border=2pt]{standalone}
\usepackage{tikz}
% =====================================================================
%  Shared TikZ style for the BISCCITS SNN workshop figures.
%  Every figure file uses:
%      \documentclass[border=2pt]{standalone}
%      \usepackage{tikz}
%      % =====================================================================
%  Shared TikZ style for the BISCCITS SNN workshop figures.
%  Every figure file uses:
%      \documentclass[border=2pt]{standalone}
%      \usepackage{tikz}
%      % =====================================================================
%  Shared TikZ style for the BISCCITS SNN workshop figures.
%  Every figure file uses:
%      \documentclass[border=2pt]{standalone}
%      \usepackage{tikz}
%      \input{style.tex}
%  so that all figures share one visual language (see SKILL: tikz-figure).
% =====================================================================
\usepackage{amsmath}
\usepackage{amssymb}   % \mathbb, etc.
\usetikzlibrary{arrows.meta, positioning, calc, backgrounds, fit,
                decorations.pathmorphing, shapes.misc}

% ----- palette: matches the matplotlib tab: colours used in the notebook -----
\definecolor{tabblue}{HTML}{1F77B4}    % input / current
\definecolor{taborange}{HTML}{FF7F0E}  % accent / secondary
\definecolor{tabgreen}{HTML}{2CA02C}   % SNN / "good"
\definecolor{tabred}{HTML}{D62728}     % threshold / reset
\definecolor{tabpurple}{HTML}{9467BD}  % membrane potential V
\definecolor{tabgray}{HTML}{7F7F7F}    % leak / inactive / captions
\definecolor{inkblack}{HTML}{222222}   % spikes / outlines / text

% ----- typography -----
\tikzset{every node/.style={font=\small}}

% ----- canonical element styles -----
\tikzset{
  % the LIF neuron glyph -- reuse this EVERYWHERE a neuron appears
  lifnode/.style={circle, draw=inkblack, line width=0.7pt,
                  minimum size=11mm, fill=tabpurple!12, inner sep=1pt},
  % a smaller neuron for dense multi-neuron layouts
  lifdot/.style={circle, draw=inkblack, line width=0.6pt,
                 minimum size=7mm, fill=tabpurple!12, inner sep=0pt},
  iolabel/.style={font=\footnotesize, inner sep=1.5pt},
  current/.style={-{Stealth[length=2.2mm]}, tabblue,  line width=0.9pt}, % input current
  synapse/.style={-{Stealth[length=2.2mm]}, inkblack, line width=0.8pt}, % forward connection
  spike/.style  ={-{Stealth[length=2.4mm]}, inkblack, line width=1.0pt}, % spike output
  leak/.style   ={-{Stealth[length=1.8mm]}, tabgray,  line width=0.8pt, dashed}, % leak / recurrence
  inactive/.style={tabgray!55, line width=0.6pt},      % "skipped" / silent path
  thresholdline/.style={tabred, dashed, line width=0.8pt},
  block/.style={draw=inkblack, rounded corners=2pt, fill=white,
                minimum height=8mm, minimum width=16mm, inner sep=3pt},
  layerbox/.style={draw=tabgray, rounded corners=2pt, dashed, inner sep=3mm},
  fignote/.style={font=\footnotesize\itshape, text=tabgray},
}

% ----- canonical symbols (write them identically in every figure) -----
\newcommand{\Vmem}{V}
\newcommand{\Vthr}{V_{\mathrm{thr}}}
\newcommand{\Vreset}{V_{\mathrm{reset}}}
\newcommand{\Iin}{I}
\newcommand{\Sout}{S}
\newcommand{\bbeta}{\beta}
\newcommand{\Wlayer}{W}
\newcommand{\slope}{k}

%  so that all figures share one visual language (see SKILL: tikz-figure).
% =====================================================================
\usepackage{amsmath}
\usepackage{amssymb}   % \mathbb, etc.
\usetikzlibrary{arrows.meta, positioning, calc, backgrounds, fit,
                decorations.pathmorphing, shapes.misc}

% ----- palette: matches the matplotlib tab: colours used in the notebook -----
\definecolor{tabblue}{HTML}{1F77B4}    % input / current
\definecolor{taborange}{HTML}{FF7F0E}  % accent / secondary
\definecolor{tabgreen}{HTML}{2CA02C}   % SNN / "good"
\definecolor{tabred}{HTML}{D62728}     % threshold / reset
\definecolor{tabpurple}{HTML}{9467BD}  % membrane potential V
\definecolor{tabgray}{HTML}{7F7F7F}    % leak / inactive / captions
\definecolor{inkblack}{HTML}{222222}   % spikes / outlines / text

% ----- typography -----
\tikzset{every node/.style={font=\small}}

% ----- canonical element styles -----
\tikzset{
  % the LIF neuron glyph -- reuse this EVERYWHERE a neuron appears
  lifnode/.style={circle, draw=inkblack, line width=0.7pt,
                  minimum size=11mm, fill=tabpurple!12, inner sep=1pt},
  % a smaller neuron for dense multi-neuron layouts
  lifdot/.style={circle, draw=inkblack, line width=0.6pt,
                 minimum size=7mm, fill=tabpurple!12, inner sep=0pt},
  iolabel/.style={font=\footnotesize, inner sep=1.5pt},
  current/.style={-{Stealth[length=2.2mm]}, tabblue,  line width=0.9pt}, % input current
  synapse/.style={-{Stealth[length=2.2mm]}, inkblack, line width=0.8pt}, % forward connection
  spike/.style  ={-{Stealth[length=2.4mm]}, inkblack, line width=1.0pt}, % spike output
  leak/.style   ={-{Stealth[length=1.8mm]}, tabgray,  line width=0.8pt, dashed}, % leak / recurrence
  inactive/.style={tabgray!55, line width=0.6pt},      % "skipped" / silent path
  thresholdline/.style={tabred, dashed, line width=0.8pt},
  block/.style={draw=inkblack, rounded corners=2pt, fill=white,
                minimum height=8mm, minimum width=16mm, inner sep=3pt},
  layerbox/.style={draw=tabgray, rounded corners=2pt, dashed, inner sep=3mm},
  fignote/.style={font=\footnotesize\itshape, text=tabgray},
}

% ----- canonical symbols (write them identically in every figure) -----
\newcommand{\Vmem}{V}
\newcommand{\Vthr}{V_{\mathrm{thr}}}
\newcommand{\Vreset}{V_{\mathrm{reset}}}
\newcommand{\Iin}{I}
\newcommand{\Sout}{S}
\newcommand{\bbeta}{\beta}
\newcommand{\Wlayer}{W}
\newcommand{\slope}{k}

%  so that all figures share one visual language (see SKILL: tikz-figure).
% =====================================================================
\usepackage{amsmath}
\usepackage{amssymb}   % \mathbb, etc.
\usetikzlibrary{arrows.meta, positioning, calc, backgrounds, fit,
                decorations.pathmorphing, shapes.misc}

% ----- palette: matches the matplotlib tab: colours used in the notebook -----
\definecolor{tabblue}{HTML}{1F77B4}    % input / current
\definecolor{taborange}{HTML}{FF7F0E}  % accent / secondary
\definecolor{tabgreen}{HTML}{2CA02C}   % SNN / "good"
\definecolor{tabred}{HTML}{D62728}     % threshold / reset
\definecolor{tabpurple}{HTML}{9467BD}  % membrane potential V
\definecolor{tabgray}{HTML}{7F7F7F}    % leak / inactive / captions
\definecolor{inkblack}{HTML}{222222}   % spikes / outlines / text

% ----- typography -----
\tikzset{every node/.style={font=\small}}

% ----- canonical element styles -----
\tikzset{
  % the LIF neuron glyph -- reuse this EVERYWHERE a neuron appears
  lifnode/.style={circle, draw=inkblack, line width=0.7pt,
                  minimum size=11mm, fill=tabpurple!12, inner sep=1pt},
  % a smaller neuron for dense multi-neuron layouts
  lifdot/.style={circle, draw=inkblack, line width=0.6pt,
                 minimum size=7mm, fill=tabpurple!12, inner sep=0pt},
  iolabel/.style={font=\footnotesize, inner sep=1.5pt},
  current/.style={-{Stealth[length=2.2mm]}, tabblue,  line width=0.9pt}, % input current
  synapse/.style={-{Stealth[length=2.2mm]}, inkblack, line width=0.8pt}, % forward connection
  spike/.style  ={-{Stealth[length=2.4mm]}, inkblack, line width=1.0pt}, % spike output
  leak/.style   ={-{Stealth[length=1.8mm]}, tabgray,  line width=0.8pt, dashed}, % leak / recurrence
  inactive/.style={tabgray!55, line width=0.6pt},      % "skipped" / silent path
  thresholdline/.style={tabred, dashed, line width=0.8pt},
  block/.style={draw=inkblack, rounded corners=2pt, fill=white,
                minimum height=8mm, minimum width=16mm, inner sep=3pt},
  layerbox/.style={draw=tabgray, rounded corners=2pt, dashed, inner sep=3mm},
  fignote/.style={font=\footnotesize\itshape, text=tabgray},
}

% ----- canonical symbols (write them identically in every figure) -----
\newcommand{\Vmem}{V}
\newcommand{\Vthr}{V_{\mathrm{thr}}}
\newcommand{\Vreset}{V_{\mathrm{reset}}}
\newcommand{\Iin}{I}
\newcommand{\Sout}{S}
\newcommand{\bbeta}{\beta}
\newcommand{\Wlayer}{W}
\newcommand{\slope}{k}


\begin{document}
\begin{tikzpicture}[x=1cm, y=1cm]

  % =================================================================
  % Reusable shared input glyph: a small T x C grid (the same series).
  % We draw it three times (one per architecture), vertically stacked.
  % \drawinput{<x>}{<y>}  places the grid with its CENTRE near (x,y).
  % =================================================================
  \def\drawinput#1#2{%
    \begin{scope}[shift={(#1,#2)}]
      % 3 rows (channels, abbreviated) x 4 cols (time, abbreviated)
      \foreach \r in {0,1,2}{
        \foreach \c in {0,1,2,3}{
          \draw[draw=tabgray, line width=0.4pt, fill=tabblue!10]
                (\c*0.26,\r*0.26) rectangle (\c*0.26+0.26,\r*0.26+0.26);
        }
      }
      \node[iolabel, text=tabgray, scale=0.8] at (1.18,0.39) {$\cdots$};
      \node[iolabel, text=tabblue, scale=0.85,
            rotate=90, anchor=south] at (-0.12,0.39) {$x\in\mathbb{R}^{6}$};
      \node[iolabel, text=tabblue, scale=0.85] at (0.52,-0.34) {time $T$};
    \end{scope}}

  % horizontal anchors (shared logits column; mean-over-time for both recurrent rows)
  \def\colInput{0}      % left edge of input grid
  \def\xcellA{2.9}      % first unrolled recurrent cell
  \def\xcellB{4.3}      % second unrolled cell
  \def\xcellC{5.7}      % third unrolled cell
  \def\xmean{7.5}       % mean-over-time block (GRU + SNN)
  \def\colLog{9.3}      % shared logits column

  % row centres for the three architectures (top -> bottom)
  \def\rowMLP{4.6}
  \def\rowGRU{2.3}
  \def\rowSNN{0.0}

  % small helper for the "logits" output marker (a 3-bar mini-vector)
  \def\drawlogits#1#2{%
    \begin{scope}[shift={(#1,#2)}]
      \draw[draw=inkblack, line width=0.5pt] (-0.16,-0.26) rectangle (0.16,0.26);
      \draw[draw=inkblack, line width=0.35pt] (-0.16,-0.085) -- (0.16,-0.085);
      \draw[draw=inkblack, line width=0.35pt] (-0.16, 0.085) -- (0.16, 0.085);
      \node[iolabel, right=1pt] at (0.16,0) {logits};
    \end{scope}}

  % =================================================================
  % ROW 1 : MLP   (per-timestep MLP, NO hidden state) -> mean over time -> logits
  % Same layout as GRU/SNN, but with NO arrows between blocks (not recurrent).
  % =================================================================
  \node[iolabel, font=\footnotesize\bfseries] at (\colInput-0.2,\rowMLP+0.95) {MLP};
  \drawinput{\colInput}{\rowMLP-0.4}
  \node[block, minimum width=11mm] (m1) at (\xcellA,\rowMLP) {MLP};
  \node[block, minimum width=11mm] (m2) at (\xcellB,\rowMLP) {MLP};
  \node[block, minimum width=11mm] (m3) at (\xcellC,\rowMLP) {MLP};
  \draw[synapse] (\colInput+1.05,\rowMLP) -- (m1.west);
  % each MLP reads its own input column; deliberately NO arrows between blocks
  \foreach \m in {m1, m2, m3}{
    \draw[current] ($(\m.south)+(0,-0.45)$) -- (\m.south);
  }
  \node[block, minimum width=14mm, align=center]
        (mmean) at (\xmean,\rowMLP) {mean over\\ time};
  \draw[synapse] (m3.east) -- (mmean.west);
  \draw[synapse] (mmean.east) -- (\colLog-0.2,\rowMLP);
  \drawlogits{\colLog}{\rowMLP}
  \node[fignote] at (\xcellB,\rowMLP-0.95) {no hidden state passed through time};

  % =================================================================
  % ROW 2 : GRU  (unrolled cell, hidden state h_t) -> mean over time -> logits
  % =================================================================
  \node[iolabel, font=\footnotesize\bfseries] at (\colInput-0.2,\rowGRU+0.95) {GRU};
  \drawinput{\colInput}{\rowGRU-0.4}
  \node[block, minimum width=11mm] (g1) at (\xcellA,\rowGRU) {GRU};
  \node[block, minimum width=11mm] (g2) at (\xcellB,\rowGRU) {GRU};
  \node[block, minimum width=11mm] (g3) at (\xcellC,\rowGRU) {GRU};
  \draw[synapse] (\colInput+1.05,\rowGRU) -- (g1.west);
  \draw[leak] (g1.east) -- (g2.west) node[midway, above, iolabel, text=tabgray] {$h_t$};
  \draw[leak] (g2.east) -- (g3.west) node[midway, above, iolabel, text=tabgray] {$h_t$};
  \foreach \g in {g1, g2, g3}{
    \draw[current] ($(\g.south)+(0,-0.45)$) -- (\g.south);
  }
  \node[block, minimum width=14mm, align=center]
        (gmean) at (\xmean,\rowGRU) {mean over\\ time};
  \draw[synapse] (g3.east) -- (gmean.west);
  \draw[synapse] (gmean.east) -- (\colLog-0.2,\rowGRU);
  \drawlogits{\colLog}{\rowGRU}
  \node[fignote] at (\xcellB,\rowGRU-0.95) {recurrent, real-valued};

  % =================================================================
  % ROW 3 : SNN  (unrolled LIF dots, beta state, spikes) -> mean over time -> logits
  % =================================================================
  \node[iolabel, font=\footnotesize\bfseries] at (\colInput-0.2,\rowSNN+0.95) {SNN};
  \drawinput{\colInput}{\rowSNN-0.4}
  \node[lifdot] (s1) at (\xcellA,\rowSNN) {};
  \node[lifdot] (s2) at (\xcellB,\rowSNN) {};
  \node[lifdot] (s3) at (\xcellC,\rowSNN) {};
  \draw[synapse] (\colInput+1.05,\rowSNN) -- (s1.west);
  \draw[leak] (s1.east) -- (s2.west) node[midway, above, iolabel, text=tabgray] {$\bbeta$};
  \draw[leak] (s2.east) -- (s3.west) node[midway, above, iolabel, text=tabgray] {$\bbeta$};
  \foreach \s in {s1, s2, s3}{
    \draw[current] ($(\s.south)+(0,-0.45)$) -- (\s.south);
    \draw[spike] (\s.north) -- ($(\s.north)+(0,0.32)$);
  }
  \node[block, minimum width=14mm, align=center]
        (smean) at (\xmean,\rowSNN) {mean over\\ time};
  \draw[synapse] (s3.east) -- (smean.west);
  \draw[synapse] (smean.east) -- (\colLog-0.2,\rowSNN);
  \drawlogits{\colLog}{\rowSNN}
  \node[fignote] at (\xcellB,\rowSNN-0.95) {recurrent, spiking};

\end{tikzpicture}
\end{document}
